%! Modeling with Exponential Functions: Interest, Half-Life, and e — Unit 6, Lesson 6.4 — Exit Ticket (Answer Key)
\documentclass[10pt]{article}
\usepackage{algebra2-article}
\usepackage{algebra2-key}   % pulls in -boxes; do NOT also load -boxes
\newcommand{\TallMath}[1]{$\displaystyle #1\rule[-1.4em]{0pt}{3.2em}$}

\begin{document}

\pageheader{Unit 6, Lesson 6.4}{Exit Ticket --- Answer Key}
\namedateperiod

\vspace{0.08in}

No notes. On items 1 and 2, write the substitution before you write the answer --- the model \emph{is}
part of the answer. Round money to the nearest cent.

\begin{enumerate}[leftmargin=1.8em, itemsep=11pt]
  \item \textbf{Compound interest.} \$4000 is invested at $3\%$ compounded \textbf{monthly} for
        $5$ years.

        \vspace{4pt}
        Base $1+\frac{r}{n}=$ \ans{$1.0025$} \qquad Exponent $nt=$ \ans{$60$}

        \vspace{5pt}
        $A=$ \ans{$4000(1.0025)^{60}$} $=$ \ans{\$4646.47}

  \item \textbf{Half-life.} A $50$ mg dose of a drug has a half-life of $4$ days.

        \vspace{4pt}
        $A(t)=$ \ans{$50\left(\frac12\right)^{t/4}$} \hfill Amount left after $12$ days: \ans{6.25 mg}

        \vspace{5pt}
        Will the amount in the patient's body ever reach exactly $0$ mg according to this model?
        Circle one --- \textbf{yes} \; / \; \textcolor{keyred}{\textbf{no}} --- and name the feature of the graph that
        settles it. \ansline{The horizontal asymptote at $A=0$: the range is $(0,50]$, so the curve approaches zero forever without touching it.}

  \item \textbf{Multiple choice.} An investment of \$1500 earns $4.8\%$ annual interest compounded
        \textbf{quarterly}. Which expression gives its value after $7$ years?
        \begin{enumerate}[label=\Alph*., leftmargin=1.9em, itemsep=1pt]
          \item $1500(1.048)^{7}$
          \item $1500(1.012)^{7}$
          \item $1500(1.048)^{28}$
          \item $1500(1.012)^{28}$ \quad \textcolor{keyred}{\textbf{$\leftarrow$ correct}}
        \end{enumerate}

  \item \textbf{The wall.} To find when that \$1500 doubles, you would solve
        $1500(1.012)^{4t}=3000$, which becomes $(1.012)^{4t}=2$. Yesterday's common-base method cannot
        finish this. Does the equation still \emph{have} an answer? Circle one ---
        \textcolor{keyred}{\textbf{yes}} \; /
        \; \textbf{no} --- and say how you know. \ansline{$2$ is positive, so it is in the range of the growth curve --- the line $y=2$ really does cross it. Only the method fails: $2$ is not a power of $1.012$.}
\end{enumerate}

\begin{teachernote}
\textbf{Answer 3: D.} The three distractors are the lesson's three failure modes, and each one names a
different re-teach. \textbf{A} splits neither the rate nor the time --- the student is still in Lesson
6.1 and has not registered that compounding happens more than once a year. \textbf{C} divides the time
but not the rate (Renata's error from the Activity); \textbf{B} divides the rate but not the time
(Miles's). Sort by distractor, not by right/wrong: A needs the Hook again, B and C each need one line
of the compound-interest box.

\textbf{Item 4 is the most informative line on the page and should be read first.} ``No'' is wrong, and
it is the misconception Lesson 6.3's exit ticket was written to catch, arriving one day later in a real
context. A class with several ``no'' answers cannot start Unit 7 --- five minutes on the Section 4
graph, with the gold line visibly crossing the curve, fixes it. Students who write ``yes, it's about
$14$ or $15$ years'' have transferred the graph reading and should be told so.

\textbf{Item 2's last part is Lesson 6.0 arriving in a context where it matters clinically.} Accept
``horizontal asymptote'' or ``the range never includes $0$''; do not accept ``because you can always
halve again'' without the graph feature named. A student who answers ``yes'' has reverted to the Unit 5
endpoint habit that 6.0 was built to break.

\textbf{Sort the collected tickets into four piles:} \emph{got it} / \emph{setup} (item 1 or 3 wrong ---
base and exponent disagree) / \emph{asymptote} (item 2's last part) / \emph{the wall} (item 4 answered
``no''). The \emph{setup} pile is the largest on a normal day and needs practice, not re-teaching. The
\emph{wall} pile is the one that changes how Lesson 6.5 and Unit 7 open.
\end{teachernote}

\end{document}
